\section{Flat-Stable LoRA (FS-LoRA)}
\label{sec:method}

We present FS-LoRA as the solution to a constrained optimization problem:
\emph{minimize current-task sharpness subject to a Fisher-weighted drift budget}.
This framing motivates all design elements---decoupled gradient combination,
trace-normalized Fisher, and dual ascent---as components of a principled Lagrangian
optimizer rather than a heuristic sum.

\subsection{Problem Formulation: Constrained Flat-Stable Optimization}
\label{sec:lagrangian}

The joint goal of flatness and stability can be stated as a constrained program:
\begin{equation}
  \min_\phi \;\Ltask(\phi) + \lambdaflat S^{(1)}_\Delta(\phi)
  \quad\text{s.t.}\quad
  \Dtilde^{\mathrm{Fisher}}(\phi) \;\leq\; \delta,
  \label{eq:constrained}
\end{equation}
where $S^{(1)}_\Delta(\phi) = \|\ggam - \gclean\|$ is the AS$^1$ sharpness
(gradient norm at the perturbed point, measuring first-order curvature),
and $\Dtilde^{\mathrm{Fisher}}(\phi) = \sum_{l,t'} \Ftilde_{t',l} \cdot \|\DeltaW_l - \DeltaW^*_{t',l}\|_F^2 / \|\Delta\phi\|^2$
is the normalized Fisher-weighted drift measuring how far the adapter has moved
along past-task sensitive directions.
The constraint $\Dtilde^{\mathrm{Fisher}} \leq \delta$ directly encodes the stability requirement:
the adapter update must not excessively disturb old-task sensitive directions.

\textbf{Penalty method approximation.}
Introducing the Lagrangian dual variable $\lambdaewc \geq 0$ for the constraint
yields the augmented objective:
\begin{equation}
  \min_\phi \;\Ltask(\phi) + \lambdaflat S^{(1)}_\Delta(\phi)
  + \tfrac{\lambdaewc}{2}\,\Lfisher(\phi),
  \label{eq:lagrangian}
\end{equation}
where $\Lfisher(\phi) = \sum_{l,t'} (\Ftilde_{t',l}+\eta)\|\DeltaW_l-\DeltaW^*_{t',l}\|_F^2$
is the trace-normalized EWC penalty (Section~\ref{sec:trace_norm}).
$\lambdaewc$ is the Lagrangian \emph{dual estimate}: it controls how strongly the constraint
is penalized.
Unlike treating $\lambdaewc$ as a static hyperparameter, the Lagrangian perspective
motivates \emph{adapting} it to the observed constraint violation---this is dual ascent
(Section~\ref{sec:dual_ascent}).

\textbf{Why this framing matters.}
The formulation~\eqref{eq:constrained}--\eqref{eq:lagrangian} is not merely aesthetic.
It makes precise that:
(1)~the Fisher term is a \emph{constraint} on drift direction, not an independent loss to minimize;
(2)~the flatness term shapes the geometry of the primal update, independent of the constraint;
(3)~the near-orthogonality between the flatness and Fisher gradients
(Section~\ref{sec:orthogonality}) justifies treating these as geometrically independent channels---
the primal update does not inadvertently violate the constraint, and the constraint does not
distort the flatness direction.

\subsection{The Coupled-Gradient Problem}

A natural baseline is to apply a sharpness-aware optimizer to the combined objective:
\begin{equation}
  \min_\phi \; \Ltask(\phi) + \lambdaewc \Lfisher(\phi).
  \label{eq:coupled}
\end{equation}
This means the GAM perturbation direction is:
\begin{equation}
  \hat{\phi}_{\mathrm{coupled}} = \phi + \rho \cdot \frac{\nabla_\phi (\Ltask + \lambdaewc \Lfisher)}
  {\|\nabla_\phi (\Ltask + \lambdaewc \Lfisher)\|},
\end{equation}
so the flatness perturbation is jointly influenced by both the current-task loss and the Fisher penalty.
This coupling has three problems:
(1)~the flatness component no longer measures only current-task sharpness, making mechanistic
interpretation ambiguous;
(2)~the Fisher penalty gradient, which has norm $\mathcal{O}(10^{-6})$ in raw scale,
may have negligible influence on the perturbation direction, making the combined optimization
effectively equivalent to GAM on $\Ltask$ alone;
(3)~\emph{crucially, coupling destroys the near-orthogonality between the two objectives.}
When the GAM closure includes $\Lfisher$, the flat gradient gains a component
proportional to $J^\top F_{t'} J\epsilon$ (see Section~\ref{sec:orthogonality}),
which projects onto the Fisher eigenspace and produces a non-zero cosine.
The empirical gap between Layer~1 (coupled, FAA$=73.23$) and Layer~2 (decoupled, FAA$=73.85$)
is a direct consequence of this distortion.

\subsection{Decoupled Gradient Combination}
\label{sec:decoupled}

FS-LoRA decouples the three gradient sources explicitly.
Keeping the GAM closure strictly on $\Ltask$ ensures that the flatness component
does not acquire a Fisher-aligned term $J^\top F_{t'} J\epsilon$
(see Section~\ref{sec:orthogonality}).
The three components are:
\begin{align}
  \gclean &= \nabla_\phi \Ltask(\phi), \label{eq:gclean}\\
  \gflat  &= \nabla_\phi \Ltask\bigl(\phi + \rho \cdot \gclean/\|\gclean\|\bigr)
             - \gclean, \label{eq:gflat}\\
  \gfisher &= \nabla_\phi \Lfisher(\phi). \label{eq:gfisher}
\end{align}

The near-orthogonality between $\gflat$ and $\gfisher$ (Section~\ref{sec:orthogonality})
justifies treating these as geometrically independent channels: the flatness update does not
inadvertently push the adapter toward Fisher-sensitive directions.
We monitor the per-step cosine $\cosff = \gflat^\top\gfisher / (\|\gflat\|\,\|\gfisher\|+\epsilon)$
as an empirical diagnostic; across all experiments it remains in $[-0.03, -0.001]$,
confirming that the decoupled combination is well-conditioned.
The final update is:
\begin{equation}
  \gupdate = \gclean + \lambdaflat\,\gflat + \gfisher.
  \label{eq:update}
\end{equation}

\subsection{Trace-Normalized Fisher Penalty}
\label{sec:trace_norm}

We use the trace-normalized Fisher (Eq.~\eqref{eq:trace_norm}) in the EWC penalty:
\begin{equation}
  \Lfisher(\phi) = \frac{\lambdaewc}{2} \sum_{l,t'<t}
  \bigl(\Ftilde_{t',l} + \eta\bigr) \cdot \bigl\|\DeltaW_l - \DeltaW^*_{t',l}\bigr\|_F^2.
\end{equation}
The Fisher gradient with respect to adapter factors is derived via the chain rule
through $\DeltaW_l = B_l A_l$:
\begin{align}
  \frac{\partial \Lfisher}{\partial B_l} &= \lambdaewc \sum_{t'} (\Ftilde_{t',l}+\eta)
    (\DeltaW_l - \DeltaW^*_{t',l}) A_l^\top, \\
  \frac{\partial \Lfisher}{\partial A_l} &= \lambdaewc \sum_{t'} (\Ftilde_{t',l}+\eta)
    B_l^\top (\DeltaW_l - \DeltaW^*_{t',l}).
\end{align}

Table~\ref{tab:norm_ratio} shows the effect of trace normalization: without it,
the Fisher-to-clean gradient norm ratio is $0.003$ at $\lambdaewc=20$,
meaning the Fisher penalty is structurally inactive.
After normalization, even with $\lambdaewc=2000$ the ratio reaches $0.025$--$0.19$,
showing that the Fisher component is active and makes a genuine contribution to the update.

\begin{table}[h]
\caption{Effect of trace normalization on Fisher activity (CUB200, seed 0).
Fisher/clean ratio = $\|\gfisher\| / \|\gclean\|$ (mean over tasks).
Without normalization, even at $\lambdaewc=1000$ the ratio reaches only $0.06$;
with normalization and $\lambdaewc=2000$ it is $0.025$, using a dimensionless hyperparameter.}
\label{tab:norm_ratio}
\centering
\begin{tabular}{lcccc}
\toprule
Method & $\lambdaewc$ & Normalized & Fisher/clean ratio & CNN FAA \\
\midrule
Layer 2 (raw) & 20   & No  & 0.0029 & 73.845 \\
Layer 2 (raw) & 100  & No  & 0.0137 & 74.182 \\
Layer 2 (raw) & 1000 & No  & 0.0588 & 74.220 \\
\midrule
Layer 3 (ours) & 2000 & Yes & 0.0248 & 74.224 \\
Layer 3 (ours) & 5000 & Yes & ---    & 74.181 \\
\bottomrule
\end{tabular}
\end{table}

\subsection{Fisher Accumulation with Exponential Decay}
\label{sec:fisher_accum}

After completing task $t$, we compute the empirical Fisher over a mini-batch of task-$t$ data
and accumulate with decay rate $\gamma$:
\begin{equation}
  F_t \leftarrow \gamma \cdot F_{t-1} + (1-\gamma) \cdot \hat{F}_t,
\end{equation}
where $\hat{F}_t$ is the mini-batch Fisher estimate.
We use $\gamma = 0.9$ throughout.
The decay prevents older task Fishers from dominating the penalty as the task count grows.

\subsection{Dual Ascent for Adaptive Constraint Enforcement}
\label{sec:dual_ascent}

The Lagrangian dual variable $\lambdaewc$ in~\eqref{eq:lagrangian} should adapt to
the observed constraint violation rather than remain fixed.
The standard dual ascent update is:
\begin{equation}
  \lambdaewc \;\leftarrow\; \max\!\bigl(\lambdaewc^{\min},\;
    \min\!\bigl(\lambdaewc^{\max},\;
      \lambdaewc + \alpha_{\mathrm{dual}}\,
      \bigl(\Dtilde^{\mathrm{Fisher}}(\phi) - \delta\bigr)\bigr)\bigr),
  \label{eq:dual_ascent}
\end{equation}
applied once per task boundary after computing the post-task Fisher drift
$\Dtilde^{\mathrm{Fisher}}$.
When the drift exceeds $\delta$, the penalty is strengthened; when the constraint
is comfortably satisfied, $\lambdaewc$ can decay toward $\lambdaewc^{\min}$.
This closes the outer loop of the Lagrangian: the inner loop (per step) optimizes
the primal variables $\phi$ at fixed $\lambdaewc$; the outer loop (per task) adjusts
$\lambdaewc$ to steer toward constraint satisfaction.

\textbf{Inactive in current experiments.}
On all benchmarks in Section~\ref{sec:experiments}, the normalized Fisher drift
$\Dtilde^{\mathrm{Fisher}}$ remained below the threshold $\delta$ throughout training,
so~\eqref{eq:dual_ascent} did not engage and $\lambdaewc$ was held fixed.
This is consistent with the orthogonality finding: because $\gflat$ and $\gfisher$ are
near-orthogonal, the flatness update does not push the adapter toward Fisher-sensitive
directions, keeping drift naturally low.
Dual ascent is expected to become active on task sequences with high semantic overlap,
where the Fisher drift constraint is binding.
We include it as the principled outer-loop completion of the Lagrangian design.

\subsection{Algorithm Summary}
\label{sec:algorithm}

\begin{algorithm}[H]
\caption{FS-LoRA Training}
\label{alg:fslora}
\begin{algorithmic}[1]
\Require Tasks $\mathcal{T}_1, \ldots, \mathcal{T}_T$; hyperparameters $\lambdaflat, \lambdaewc, \rho, \gamma, \delta, \alpha_{\mathrm{dual}}$
\State Initialize adapter $\phi$, Fisher $\{F_0^l = 0\}$, snapshots $\{\DeltaW^*_0\}$
\For{task $t = 1, \ldots, T$}
  \For{each mini-batch $\mathcal{B} \subset \mathcal{T}_t$}
    \State $\gclean \leftarrow \nabla_\phi \Ltask(\phi; \mathcal{B})$ \Comment{clean gradient}
    \State $\gflat \leftarrow \nabla_\phi \Ltask(\phi + \rho\,\gclean/\|\gclean\|; \mathcal{B}) - \gclean$ \Comment{AS$^1$ flatness component}
    \State $\gfisher \leftarrow \nabla_\phi \Lfisher(\phi)$ using stored $\{F_{t'}, \DeltaW^*_{t'}\}_{t'<t}$
    \State $\phi \leftarrow \phi - \eta_{\mathrm{lr}} \cdot (\gclean + \lambdaflat\,\gflat + \gfisher)$
  \EndFor
  \State Snapshot: $\DeltaW^*_t \leftarrow \{B_l A_l\}$ for all $l$
  \State Compute $\hat{F}_t$ from a mini-batch of $\mathcal{T}_t$; normalize by trace
  \State Accumulate: $F_t \leftarrow \gamma F_{t-1} + (1-\gamma)\hat{F}_t$
  \If{$t > 1$} \Comment{dual ascent outer loop; inactive when $\Dtilde^{\mathrm{Fisher}} \leq \delta$}
    \State $\lambdaewc \leftarrow \mathrm{clip}\!\bigl(\lambdaewc + \alpha_{\mathrm{dual}}(\Dtilde^{\mathrm{Fisher}} - \delta),\;\lambdaewc^{\min},\;\lambdaewc^{\max}\bigr)$
  \EndIf
\EndFor
\end{algorithmic}
\end{algorithm}

The algorithm implements the two-level Lagrangian structure from
Section~\ref{sec:lagrangian}: (1)~per-step decoupled gradient combination
(always active); (2)~per-task dual ascent on $\lambdaewc$
(inactive when $\Dtilde^{\mathrm{Fisher}}\leq\delta$).
The computational overhead relative to standard SGD is two forward passes per batch
(for $\gclean$ and $\ggam$) plus a Fisher gradient computation.
The Fisher accumulation and drift evaluation at task boundaries require a single pass
over a held-out mini-batch (we use up to 100 batches, controlled by
\texttt{ewc\_max\_batches}).
